\chapter{TDD - Test Driven Development}
\section{Definizione}
È una pratica che \uline{prevede di scrivere il test unitario prima del codice di produzione}.

Il programmatore scrive il test immaginando che \textbf{l'interfaccia esista già}, e solo successivamente scrive il codice per farlo passare.

\section{Cosa è un test unitario?}
Un test unitario è un test che si concentra sulla verifica di singole classi e metodi all'interno del codice.

\subsection{Schema di test}
\begin{enumerate}
    \item \textbf{Preparazione:}\\ Si crea l'oggetto da verificare, chiamato anche fixture
    \item \textbf{Esecuzione:}\\ Si eseguono specifiche azioni sulla fixture per attivare le operazioni che si desidera testare
    \item \textbf{Verifica:}\\ Si controlla che i risultati ottenuti dal codice corrispondano ai risultati previsti
    \item \textbf{Rilascio:}\\ Una fase opzionale in cui si ripuliscono o si rilasciono gli oggetti e le risorse che sono state utilizzate durante il test
\end{enumerate}

\subsection{Tipi di test}
\begin{itemize}
    \item Test di unità
    \item Test di integrazione
    \item Test di sistema
    \item Test di accettazione
\end{itemize}

\section{Vantaggi}
\begin{itemize}
    \item \textbf{I test vengono effettivamente scritti:}\\ Obbligando a scrivere il test prima del codice, si ha la garanzia che i test non vengano tralasciati
    \item \textbf{Maggiore coerenza:}\\ La soddisfazione del programmatore nel far passare i test porta a una scrittura dei test più coerente e disciplinata
    \item \textbf{Chiarezza progettuale:}\\ Scrivere il test immaginando l'uso dell'oggetto aiuta a chiarire i dettagli dell'interfaccia e del comportamento atteso prima ancora di implementarli
    \item \textbf{Verifica continua:}\\ Si ottiene un sistema di verifica automatica, che è costantemente ripetibile e dimostrabile in qualsiasi momento
    \item \textbf{Fiducia nel cambiamento:}\\ La presenza di test automatici infonde fiduca nel programmatore quando è necessario modificare o fare refactoring del codice sapendo che eventuali errori verrano subito rilevati
\end{itemize}

\section{Svantaggi}
\begin{itemize}
    \item \textbf{Complessità di applicazione}
    \item \textbf{Curva di apprendimento ripida}
    \item \textbf{Maggiore tempo di sviluppo iniziale}
    \item \textbf{Aumento del codice da mantenere}
\end{itemize}

\section{Progettazione degli input}
Per testare in modo efficiente, gli input vanno raggruppati in "partizioni di equivalenza" (insiemi di dati trattati allo stesso modo dal codice). 

È fondamentale testare sia input validi che input deliberatamente sbagliati, prestando particolare attenzione ai valori ai limiti (o di frontiera), forzando errori o testando valori critici come puntatori nulli e zeri.
